% Image credits for the photographs and AI illustrations of Book 5 (University, Year 3), French edition.
% Machine-readable license records live in images/book5/CREDITS.md; keep this
% file in sync with frontmatter/image-credits-book5.tex (one \item per item there).
% Every credit, author, dataset and licence is kept exactly as in the English file.
\chapter*{Crédits des images}

Les dessins, cartes et graphiques de ce livre sont originaux~; les cartes
et graphiques tracés à partir de données publiées en citent la source dans
la légende ou dans cette liste. Les photographies sont utilisées sous leurs
licences libres respectives, avec nos remerciements à leurs auteurs~; les
liens complets vers les sources et les adresses des licences de chaque
photographie figurent dans \texttt{images/book5/CREDITS.md}, dans le dépôt
du livre.

\bigskip
À côté des photographies et des dessins originaux, certaines figures sont
des illustrations générées par IA, produites avec le générateur d'images
d'OpenAI à partir de consignes rédigées par les éditeurs du livre, et
vérifiées pour leur exactitude scientifique avant d'être retenues. Elles
ne montrent que des scènes génériques, jamais un lieu réel, une personne
réelle ou un spécimen réel. La consigne complète de chaque image générée
est consignée dans \texttt{images/book5/ai/PROMPTS.md}, dans le dépôt.

\bigskip
\noindent Photographies et données, par chapitre~:
\begin{itemize}\small
  \item Chap.~1 (sismologie et imagerie de la Terre)~: PREM, A.~M.~Dziewonski et D.~L.~Anderson (1981), tel que tabulé dans ObsPy~; ak135, B.~L.~N.~Kennett, E.~R.~Engdahl et R.~Buland (1995).
  \item Chap.~3 (la pesanteur et la figure de la Terre)~: les satellites jumeaux GRACE Follow-On, vue d'artiste, NASA/JPL-Caltech, domaine public (Wikimedia Commons). Bilan de masse des calottes glaciaires, IMBIE Team (2021), UK Polar Data Centre, Open Government Licence v3.0.
  \item Chap.~4 (géomagnétisme et paléomagnétisme)~: champ géomagnétique international de référence, quatorzième génération, IAGA Division~V Working Group V-MOD, coefficients diffusés par NOAA NCEI.
  \item Chap.~5 (modèles thermiques de la lithosphère)~: âge des fonds océaniques, M.~Seton et al. (2020), EarthByte, CC BY 4.0~; relief mondial ETOPO1, NOAA NCEI, domaine public.
  \item Chap.~6 (convection du manteau et forces motrices des plaques)~: les plages soulevées de Ballyhillin, Malin Head, A.~F.~Borchert, CC BY-SA 4.0 (Wikimedia Commons)~; PREM, A.~M.~Dziewonski et D.~L.~Anderson (1981).
  \item Chap.~7 (mécanismes de déformation et zones de cisaillement)~: mylonite œillée, Mikenorton, CC BY-SA 4.0 (Wikimedia Commons).
  \item Chap.~9 (tectonique extensive et décrochante)~: relief de type «~basin and range~» au sud de la vallée de la Mort, Ken Lund, CC BY-SA 2.0 (Wikimedia Commons).
  \item Chap.~14 (géochimie des isotopes stables)~: compilation benthique $\delta^{18}$O LR04, L.~E.~Lisiecki et M.~E.~Raymo (2005).
  \item Chap.~15 (interactions eau--roche)~: terrasses de travertin de Palette Spring, Mammoth Hot Springs, Jim Peaco, National Park Service, domaine public (Wikimedia Commons). Constantes d'équilibre de la base de données phreeqc.dat de PHREEQC, D.~L.~Parkhurst et C.~A.~J.~Appelo, U.S. Geological Survey, domaine public.
  \item Chap.~16 (stratigraphie séquentielle)~: les Book Cliffs au-dessus de Helper, Utah, David Mackie, domaine public (Wikimedia Commons)~; compilation benthique $\delta^{18}$O LR04, L.~E.~Lisiecki et M.~E.~Raymo (2005).
  \item Chap.~17 (modèles de faciès)~: stratification oblique en mamelons de la formation d'Arisaig, Michael C.~Rygel, CC BY-SA 3.0 (Wikimedia Commons), recadrée.
  \item Chap.~18 (géomorphologie quantitative)~: la gorge du Niagara et les chutes du Fer-à-Cheval, Ken Lund, CC BY-SA 2.0 (Wikimedia Commons), recadrée.
  \item Chap.~19 (dynamique des fluides géophysiques)~: vents en moyenne zonale de la réanalyse NCEP--NCAR Reanalysis~1, données fournies par NOAA PSL, Boulder, Colorado, USA~; World Ocean Atlas 2023, NOAA NCEI.
  \item Chap.~20 (dynamique de l'atmosphère)~: une dépression au large de l'Islande, NASA/GSFC, MODIS Rapid Response Team, Jacques Descloitres, domaine public (Wikimedia Commons)~; vents en moyenne zonale de la réanalyse NCEP--NCAR Reanalysis~1, données fournies par NOAA PSL, Boulder, Colorado, USA.
  \item Chap.~21 (dynamique de l'océan)~: un flotteur Argo mis à l'eau depuis un navire de recherche, NOAA Global Ocean Monitoring and Observing Program, domaine public (Wikimedia Commons).
  \item Chap.~22 (modèles de climat)~: superordinateurs du NASA Center for Climate Simulation, Goddard Space Flight Center, 2010, domaine public (Wikimedia Commons)~; température de l'air en surface et rayonnement au sommet de l'atmosphère de la réanalyse NCEP--NCAR Reanalysis~1, données fournies par NOAA PSL, Boulder, Colorado, USA.
  \item Chap.~24 (glaciologie)~: un glacier émissaire du sud-ouest du Groenland, NASA / Jim Yungel, Operation IceBridge, 2014, domaine public (Wikimedia Commons)~; altitudes de la surface du Groenland d'après ETOPO1 Ice Surface, NOAA NCEI~; étendues de la banquise arctique d'après le NSIDC Sea Ice Index, version~4 (Fetterer et al., 2025, doi:10.7265/a98x-0f50).
  \item Chap.~25 (télédétection et SIG)~: l'interférogramme 1995--2001 du centre volcanique des Three Sisters, Oregon, U.S. Geological Survey (C.~Wicks et al., 2002), domaine public (Wikimedia Commons)~; spectres de laboratoire de l'ECOSTRESS Spectral Library, version~1.0 (S.~K.~Meerdink, S.~J.~Hook, D.~A.~Roberts et E.~A.~Abbott, 2019~; A.~M.~Baldridge et al., 2009), reproduits avec l'aimable autorisation du Jet Propulsion Laboratory, California Institute of Technology, Pasadena, Californie.
  \item Chap.~26 (intérieurs et surfaces des planètes)~: Barringer Meteor Crater, Arizona, USGS/D.~Roddy, domaine public (Wikimedia Commons)~; Olympus Mons, mosaïque de l'orbiteur Viking, NASA/JPL/USGS, domaine public (Wikimedia Commons). Fonctions de production et de chronologie des cratères telles qu'elles figurent dans le logiciel Craterstats de G.~G.~Michael, d'après G.~Neukum, B.~A.~Ivanov et W.~K.~Hartmann (2001).
  \item Chap.~27 (atmosphères planétaires et habitabilité)~: un réseau de vallées asséchées sur Mars, Mars Express HRSC, ESA/DLR/FU Berlin, CC BY-SA 3.0 IGO (Wikimedia Commons)~; Ligeia Mare, Titan, mosaïque radar de Cassini, NASA/JPL/USGS, domaine public (Wikimedia Commons). Ajustements de la zone habitable de R.~K.~Kopparapu et al. (2013)~; pression de vapeur du méthane d'après le NIST Chemistry WebBook.
  \item Chap.~28 (exoplanètes)~: HIP~65426~b vue par le télescope spatial James Webb, NASA/ESA/CSA, A~Carter (UCSC), the ERS~1386 team, et A.~Pagan (STScI), CC BY~4.0 (Wikimedia Commons). Planètes confirmées d'après le NASA Exoplanet Archive, table Planetary Systems Composite Parameters, consultée le 23 octobre 2026~: cette recherche a utilisé le NASA Exoplanet Archive, exploité par le California Institute of Technology sous contrat avec la National Aeronautics and Space Administration dans le cadre de l'Exoplanet Exploration Program. Courbes masse--rayon modélisées de L.~Zeng et al. (2016, 2019).
  \item Chap.~29 (la Voie lactée et les galaxies)~: le panorama de la Voie lactée du projet GigaGalaxy Zoom, ESO/S.~Brunier, CC BY~4.0 (Wikimedia Commons)~; la galaxie spirale M101, Agence spatiale européenne et NASA (Hubble, heic0602~; données Hubble K.~D.~Kuntz, F.~Bresolin, J.~Trauger, J.~Mould et Y.-H.~Chu~; traitement de l'image Davide De~Martin), CC BY~4.0 (Wikimedia Commons). Vitesses circulaires de la Galaxie d'après A.-C.~Eilers, D.~W.~Hogg, H.-W.~Rix et M.~K.~Ness (2019)~; modèle de masse d'E.~Pouliasis, P.~Di~Matteo et M.~Haywood (2017)~; céphéides classiques OGLE du Grand Nuage de Magellan, I.~Soszy\'nski et al. (2015), Acta Astronomica 65, 297, diffusées par le CDS, Strasbourg (VizieR J/AcA/65/297).
  \item Chap.~30 (d'où viennent les éléments)~: la nébuleuse du Crabe, NASA, ESA, J.~Hester et A.~Loll (Arizona State University), domaine public (Wikimedia Commons)~; la kilonova de GW170817 dans NGC~4993, NASA et ESA, remerciements à A.~J.~Levan, N.~R.~Tanvir, A.~Fruchter et O.~Fox (Hubble, heic1717a), CC BY~4.0 (Wikimedia Commons). Abondances solaires de M.~Asplund, A.~M.~Amarsi et N.~Grevesse (2021)~; sections efficaces de capture de neutrons d'après KADoNiS v0.3 (I.~Dillmann, R.~Plag, F.~K\"appeler et T.~Rauscher)~; périodes d'après la IAEA Live Chart of Nuclides (ENSDF)~; métallicités du relevé de Genève--Copenhague réanalysées par L.~Casagrande et al. (2011), diffusées par le CDS, Strasbourg (VizieR J/A+A/530/A138).
\end{itemize}
